\begin{problem}[Committee selection]
A committee containing $4$ men and $4$ women is to be formed from a group of $10$ men and $8$ women.

In how many different ways can the committee be formed?
\end{problem}

\begin{solution}
Choose the men and women independently:
\[
\binom{10}{4}\binom{8}{4} = 210 \times 70 = 14700
\]
Hence the number of possible committees is $\boxed{14700}$.
\end{solution}

\begin{takeaways}
\item When categories are separate, count each category and multiply.
\item Committee questions are combination questions because order is irrelevant.
\end{takeaways}

\begin{problem}[License plates and careful reading]
A license plate involves $2$ digits, followed by $2$ letters, followed by $2$ more digits, all chosen at random.
\begin{enumerate}[label=\alph*)]
\item How many different license plates are possible?
\item What is the probability that the last digit is prime?
\item Find the probability that the plate says \texttt{COP789}.
\end{enumerate}
\end{problem}

\begin{solution}
We interpret ``all chosen at random'' to mean repetition is allowed.

\textbf{a)} There are $10$ choices for each digit and $26$ choices for each letter:
\[
10 \times 10 \times 26 \times 26 \times 10 \times 10 = 10^4 \cdot 26^2 = 6760000
\]
So the number of license plates is $\boxed{6760000}$.

\textbf{b)} The last digit can be any of $0$ to $9$. The prime digits are $2,3,5,7$, so there are $4$ favourable digits out of $10$:
\[
\frac{4}{10} = \frac{2}{5}
\]
Thus the probability is $\boxed{\frac{2}{5}}$.

\textbf{c)} Under the stated format, a plate has the pattern digit-digit-letter-letter-digit-digit. The string \texttt{COP789} has three letters followed by three digits, so it cannot occur.
\[
\boxed{0}
\]
\end{solution}

\begin{takeaways}
\item Before counting, check the format carefully.
\item If a requested outcome is impossible under the stated rules, its probability is $0$.
\item When repetition is not forbidden, assume repeated digits or letters are allowed.
\end{takeaways}

\begin{problem}[A mixed combinatorics review]
\begin{enumerate}[label=\arabic*.]
\item How many different ways can $5$ letters be chosen from \texttt{URGOATED} if all $4$ vowels must be included?
\item How many different ways could the letters of \texttt{COOKED} be arranged?
\item How many $4$-digit numbers can be made from the digits $1$ to $8$ without repetition?
\end{enumerate}

Explain briefly which parts are combinations and which parts are permutations.
\end{problem}

\begin{solution}
\textbf{1.} This is a combination because order does not matter. We must take $U, O, A, E$ and then choose $1$ more letter from the remaining $4$ consonants:
\[
\binom{4}{1} = 4
\]

\textbf{2.} This is a permutation with repeated letters:
\[
\frac{6!}{2!} = 360
\]

\textbf{3.} This is a permutation without repetition:
\[
8 \times 7 \times 6 \times 5 = 1680
\]

So the answers are:
\[
\boxed{4,\ 360,\ 1680}
\]
\end{solution}

\begin{takeaways}
\item Combinations are for selections where order does not matter.
\item Permutations are for arrangements where order does matter.
\item Repeated objects in a permutation require division by factorials of repeated counts.
\end{takeaways}
